\documentclass[11pt]{article}
\usepackage[margin=2.0cm]{geometry}
\usepackage{booktabs}
\usepackage{amsmath}
\usepackage{amssymb}
\usepackage{enumitem}
\usepackage{array}
\usepackage{listings}
\usepackage{xcolor}
\usepackage[T1]{fontenc}
\usepackage{lmodern}

\setlength{\parindent}{0pt}
\setlength{\parskip}{3.6pt}
\newcommand{\code}[1]{\texttt{#1}}
\newcommand{\hr}{\vspace{2pt}\hrule\vspace{6pt}}
\renewcommand{\arraystretch}{1.12}

\definecolor{kw}{RGB}{0,90,160}
\definecolor{cmt}{RGB}{110,110,110}
\lstdefinestyle{vhdl}{
  language=VHDL, basicstyle=\ttfamily\footnotesize,
  keywordstyle=\color{kw}\bfseries, commentstyle=\color{cmt}\itshape,
  showstringspaces=false, columns=fullflexible, breaklines=true,
  frame=leftline, framesep=6pt, xleftmargin=8pt, aboveskip=6pt, belowskip=6pt
}
\lstset{style=vhdl}

\begin{document}

\begin{center}
{\Large\bfseries EE\,362 --- Computer Organization Laboratory}\\[2pt]
{\large\bfseries Lab 0 --- Solution}\\[4pt]
Fall 2026 \quad$\cdot$\quad Lab TA: Gandhi Torres
\end{center}
\hr

This is the worked solution for all four Lab 0 exercises. Use it to check your own
work and to back-fill \code{lab00.tex} in your Overleaf project. Every derivation
here is written in the exact form your reports are expected to use.

\vspace{6pt}
\hr

\section*{Exercise 1 --- Majority of three}

\subsection*{1.1 Specification}

$Y = 1$ when at least two of the three inputs $A$, $B$, $C$ are $1$.

\subsection*{1.2 Truth table --- Type E (exhaustive)}

Three binary inputs, so $2^3 = 8$ rows. Every row is listed; nothing is inferred.

\begin{center}
\begin{tabular}{c|ccc|c|l}
\toprule
\# & $A$ & $B$ & $C$ & $Y$ & minterm \\
\midrule
0 & 0 & 0 & 0 & 0 & \\
1 & 0 & 0 & 1 & 0 & \\
2 & 0 & 1 & 0 & 0 & \\
3 & 0 & 1 & 1 & \textbf{1} & $m_3 = \bar{A}BC$ \\
4 & 1 & 0 & 0 & 0 & \\
5 & 1 & 0 & 1 & \textbf{1} & $m_5 = A\bar{B}C$ \\
6 & 1 & 1 & 0 & \textbf{1} & $m_6 = AB\bar{C}$ \\
7 & 1 & 1 & 1 & \textbf{1} & $m_7 = ABC$ \\
\bottomrule
\end{tabular}
\end{center}

Four rows out of eight are $1$. That is a useful sanity check before you simplify:
majority-of-three is symmetric, so the function must be true exactly half the time.

\subsection*{1.3 Boolean derivation}

\textbf{(1) Canonical --- read straight off the table.}
\[
Y = \sum m(3,5,6,7) = \bar{A}BC + A\bar{B}C + AB\bar{C} + ABC
\]

\textbf{(2) Simplified.} The trick is that $m_7 = ABC$ is adjacent to each of the
other three minterms. By the \textbf{idempotent law} ($X = X + X + X$) we may write
it three times without changing the function:
\begin{align*}
Y &= \bar{A}BC + A\bar{B}C + AB\bar{C} + \underbrace{ABC + ABC + ABC}_{\text{idempotent law}} \\
  &= (\bar{A}BC + ABC) + (A\bar{B}C + ABC) + (AB\bar{C} + ABC) \\
  &= BC(\bar{A} + A) + AC(\bar{B} + B) + AB(\bar{C} + C) && \text{\small distributive law} \\
  &= BC \cdot 1 + AC \cdot 1 + AB \cdot 1 && \text{\small complement law} \\[2pt]
Y &= \boxed{AB + AC + BC}
\end{align*}

Four three-input terms became three two-input terms: 3 AND2 gates and 1 OR3 gate.

\textbf{(3) VHDL.}
\begin{lstlisting}
Y <= (A and B) or (A and C) or (B and C);
\end{lstlisting}

\subsection*{1.4 Code}

\begin{lstlisting}
library ieee;
use ieee.std_logic_1164.all;

entity majority3 is
    port ( A, B, C : in  std_logic;
           Y       : out std_logic );
end entity majority3;

architecture behavioral of majority3 is
begin
    Y <= (A and B) or (A and C) or (B and C);
end architecture behavioral;
\end{lstlisting}

\subsection*{1.5 Testbench --- coverage argument}

Three inputs, eight combinations, and the design is purely combinational: an
exhaustive testbench is both possible and cheap, so \textbf{exhaustive is the
correct choice}. Anything less would need an argument; this needs none.

\begin{lstlisting}
check : process
    variable expect : std_logic;
begin
    for i in 0 to 7 loop
        (A, B, C) <= std_logic_vector(to_unsigned(i, 3));
        wait for 10 ns;
        -- expected: 1 when the number of set bits >= 2
        if i = 3 or i = 5 or i = 6 or i = 7 then
            expect := '1';
        else
            expect := '0';
        end if;
        assert Y = expect
            report "row " & integer'image(i) & ": Y wrong"
            severity error;
    end loop;
    report "majority3: 8/8 checks complete" severity note;
    wait;
end process check;
\end{lstlisting}

Note the \code{report} on the last line. A testbench that prints nothing on success
is indistinguishable from a testbench that did not run.

\subsection*{1.6 Where this reappears}

Hold this result. You derive it a second time in Exercise 4 without meaning to.

\vspace{4pt}\hr

\section*{Exercise 2 --- 2-to-1 multiplexer}

\subsection*{2.1 Specification}

$Y = D_0$ when $S = 0$; \quad $Y = D_1$ when $S = 1$.

\subsection*{2.2 Truth table --- two ways}

\textbf{Type E (exhaustive).} Three inputs, eight rows:

\begin{center}
\begin{tabular}{c|ccc|c|l}
\toprule
\# & $S$ & $D_1$ & $D_0$ & $Y$ & minterm \\
\midrule
0 & 0 & 0 & 0 & 0 & \\
1 & 0 & 0 & 1 & \textbf{1} & $m_1 = \bar{S}\bar{D_1}D_0$ \\
2 & 0 & 1 & 0 & 0 & \\
3 & 0 & 1 & 1 & \textbf{1} & $m_3 = \bar{S}D_1D_0$ \\
4 & 1 & 0 & 0 & 0 & \\
5 & 1 & 0 & 1 & 0 & \\
6 & 1 & 1 & 0 & \textbf{1} & $m_6 = SD_1\bar{D_0}$ \\
7 & 1 & 1 & 1 & \textbf{1} & $m_7 = SD_1D_0$ \\
\bottomrule
\end{tabular}
\end{center}

\textbf{Type C (compressed).} The same function, written by naming a \emph{source}
instead of a value:

\begin{center}
\begin{tabular}{c|c|c|l}
\toprule
\# & $S$ & $Y$ & \\
\midrule
0 & 0 & $D_0$ & pass the low input \\
1 & 1 & $D_1$ & pass the high input \\
\bottomrule
\end{tabular}
\end{center}

Two rows say everything eight rows said. This is not laziness --- it is the only
form that survives. In Lab 3 the same multiplexer is 32 bits wide, and the
exhaustive table would have $2^{65}$ rows. \textbf{Type C is how you write a table
for a datapath.}

\subsection*{2.3 Boolean derivation}

\textbf{(1) Canonical.}
\[
Y = \sum m(1,3,6,7) = \bar{S}\bar{D_1}D_0 + \bar{S}D_1D_0 + SD_1\bar{D_0} + SD_1D_0
\]

\textbf{(2) Simplified} --- factor the two pairs, then apply the complement law:
\begin{align*}
Y &= \bar{S}D_0(\bar{D_1} + D_1) + SD_1(\bar{D_0} + D_0) && \text{\small distributive law} \\
  &= \bar{S}D_0 \cdot 1 + SD_1 \cdot 1 && \text{\small complement law} \\[2pt]
Y &= \boxed{\bar{S}\,D_0 + S\,D_1}
\end{align*}

Read the answer aloud: ``$D_0$ gets through when $S$ is low, $D_1$ gets through when
$S$ is high.'' The algebra reproduced the specification exactly. That is the check.

\textbf{(3) VHDL.}
\begin{lstlisting}
Y <= D0 when S = '0' else D1;
\end{lstlisting}

The \code{when/else} form and the Boolean form synthesise to the same logic. Write
whichever states your intent more clearly --- here, the \code{when/else}, because
the specification was written as a choice, not as a formula.

\subsection*{2.4 Testbench --- coverage argument}

Exhaustive again: 8 vectors. But note what you are really testing --- for each value
of $S$, that $Y$ tracks the selected input \emph{and ignores the other one}. Rows 2
and 5 are the important ones: in each, the unselected input is $1$ and $Y$ must
still be $0$. A mux that ORs its inputs passes rows 0, 3 and 7 and fails exactly
those two.

\vspace{4pt}\hr

\section*{Exercise 3 --- 2-to-4 decoder with enable}

\subsection*{3.1 Specification}

When $E = 1$, exactly one of $D_3 \ldots D_0$ is high --- the one selected by the
two address bits $A_1 A_0$. When $E = 0$, all outputs are low.

\subsection*{3.2 Truth table --- Type E with a don't-care row}

\begin{center}
\begin{tabular}{c|ccc|cccc|l}
\toprule
\# & $E$ & $A_1$ & $A_0$ & $D_3$ & $D_2$ & $D_1$ & $D_0$ & note \\
\midrule
0 & 0 & -- & -- & 0 & 0 & 0 & 0 & disabled; address irrelevant \\
1 & 1 & 0 & 0 & 0 & 0 & 0 & \textbf{1} & \\
2 & 1 & 0 & 1 & 0 & 0 & \textbf{1} & 0 & \\
3 & 1 & 1 & 0 & 0 & \textbf{1} & 0 & 0 & \\
4 & 1 & 1 & 1 & \textbf{1} & 0 & 0 & 0 & \\
\bottomrule
\end{tabular}
\end{center}

Five rows, not eight, because row 0 uses don't-cares to collapse four rows into one.
That is legitimate \emph{and} it is a claim: it asserts the address genuinely cannot
affect the output when disabled. Your simulation has to confirm it.

Rows 1--4 each have exactly one $1$. That pattern is called \textbf{one-hot}, and it
is the single most useful encoding in this course.

\subsection*{3.3 Boolean derivation}

One equation per output --- four outputs, four equations. Each is a single minterm
of the address, gated by the enable:

\begin{align*}
D_0 &= E \cdot \bar{A_1}\bar{A_0} \\
D_1 &= E \cdot \bar{A_1}A_0 \\
D_2 &= E \cdot A_1\bar{A_0} \\
D_3 &= E \cdot A_1A_0
\end{align*}

Nothing simplifies, and nothing should: each output is already a single product
term. Written compactly,
\[
D_k = E \cdot m_k(A_1, A_0)
\]

\textbf{This is the whole idea of a decoder: it computes all four minterms of its
address at once.} It is why a decoder is the standard way to turn an address into a
select signal --- and it is exactly the mechanism you build in Lab 6, where a 5-bit
address is decoded into 32 register write-enables.

\textbf{VHDL.}
\begin{lstlisting}
D(0) <= E and (not A1) and (not A0);
D(1) <= E and (not A1) and      A0;
D(2) <= E and      A1  and (not A0);
D(3) <= E and      A1  and      A0;
\end{lstlisting}

\subsection*{3.4 The bug in the behavioural version}

The version shown in class was:

\begin{lstlisting}
A <= A1 & A0;
with A select D <= "0001" when "00",
                   "0010" when "01",
                   "0100" when "10",
                   "1000" when "11",
                   "0000" when others;
\end{lstlisting}

\textbf{The enable is never used.} \code{E} is declared in the entity and then
ignored, so this circuit implements row~0 of the truth table incorrectly: with
$E = 0$ it still asserts an output. The \code{"0000" when others} branch looks like
it handles the disabled case, but it does not --- it handles metavalues
(\code{'U'}, \code{'X'}) on the address, which is a different thing entirely.

The structural version in the same folder \emph{does} use \code{E}. So two
architectures of nominally the same circuit disagree.

\textbf{How to catch this, in order of speed:}

\begin{enumerate}[leftmargin=2em,itemsep=1pt]
\item \textbf{Analysis \& Elaboration} --- Quartus warns that input port \code{E} is
stuck at, or never used. Read your warnings.
\item \textbf{RTL Viewer} --- \code{E} appears as a pin with nothing attached to it.
This is the single clearest picture of the bug, and it takes ten seconds.
\item \textbf{Simulation} --- only if your testbench actually drives $E = 0$. If you
tested only the enabled rows, everything passes. That is why row 0 is in the table.
\end{enumerate}

\textbf{The corrected behavioural architecture:}

\begin{lstlisting}
process (E, A1, A0)
    variable Asel : std_logic_vector(1 downto 0);
begin
    Asel := A1 & A0;
    if E = '0' then
        D <= "0000";
    else
        case Asel is
            when "00"   => D <= "0001";
            when "01"   => D <= "0010";
            when "10"   => D <= "0100";
            when "11"   => D <= "1000";
            when others => D <= "0000";
        end case;
    end if;
end process;
\end{lstlisting}

\subsection*{3.5 The lesson}

Your testbench cannot find a bug in a case you never drive, and the RTL Viewer shows
you the circuit you actually described rather than the one you meant to describe.
Both of those facts cost you nothing today. In Lab 10 --- sixteen control outputs,
one row per instruction --- they are the difference between an afternoon and a week.

\vspace{4pt}\hr

\section*{Exercise 4 --- Full adder}

\subsection*{4.1 Specification}

Add three bits: $A$, $B$, and a carry in $C_{in}$. Produce a sum bit and a carry out.
The two outputs together encode the arithmetic sum: $C_{out}\,Sum = A + B + C_{in}$
as a 2-bit number.

\subsection*{4.2 Truth table --- Type E}

\begin{center}
\begin{tabular}{c|ccc|cc|c}
\toprule
\# & $A$ & $B$ & $C_{in}$ & $Sum$ & $C_{out}$ & $A{+}B{+}C_{in}$ \\
\midrule
0 & 0 & 0 & 0 & 0 & 0 & 0 \\
1 & 0 & 0 & 1 & \textbf{1} & 0 & 1 \\
2 & 0 & 1 & 0 & \textbf{1} & 0 & 1 \\
3 & 0 & 1 & 1 & 0 & \textbf{1} & 2 \\
4 & 1 & 0 & 0 & \textbf{1} & 0 & 1 \\
5 & 1 & 0 & 1 & 0 & \textbf{1} & 2 \\
6 & 1 & 1 & 0 & 0 & \textbf{1} & 2 \\
7 & 1 & 1 & 1 & \textbf{1} & \textbf{1} & 3 \\
\bottomrule
\end{tabular}
\end{center}

The right-hand column is the check: read $C_{out}Sum$ as a two-bit binary number and
it must equal the count of ones in the input. Row 7: $C_{out}Sum = 11_2 = 3$. \checkmark

\subsection*{4.3 Derivation --- $Sum$}

\textbf{(1) Canonical.}
\[
Sum = \sum m(1,2,4,7) = \bar{A}\bar{B}C_{in} + \bar{A}B\bar{C_{in}} + A\bar{B}\bar{C_{in}} + ABC_{in}
\]

\textbf{(2) Simplified --- and here the K-map refuses.} Plot $\sum m(1,2,4,7)$ and
you get a checkerboard: no two 1-cells are adjacent, so \textbf{no grouping of size
two or larger exists}. The minimal sum-of-products form \emph{is} the canonical
form: four three-input AND gates. Sum-of-products is the wrong tool for this
function.

Factor by $C_{in}$ instead:
\begin{align*}
Sum &= C_{in}(\bar{A}\bar{B} + AB) + \bar{C_{in}}(\bar{A}B + A\bar{B}) && \text{\small distributive law} \\
    &= C_{in}\cdot\overline{(A \oplus B)} + \bar{C_{in}}\cdot(A \oplus B) && \text{\small definition of XOR and XNOR} \\[2pt]
Sum &= \boxed{A \oplus B \oplus C_{in}}
\end{align*}

The last step is the definition of XOR applied to the two operands $(A \oplus B)$ and
$C_{in}$: the expression is ``one or the other but not both''.

\textbf{The lesson:} not every function has a small sum-of-products form. Parity ---
which is what $Sum$ is --- is the classic example. When a K-map gives you a
checkerboard, stop trying to group and reach for XOR.

\textbf{(3) VHDL.}
\begin{lstlisting}
Sum <= A xor B xor Cin;
\end{lstlisting}

\subsection*{4.4 Derivation --- $C_{out}$}

\textbf{(1) Canonical.}
\[
C_{out} = \sum m(3,5,6,7) = \bar{A}BC_{in} + A\bar{B}C_{in} + AB\bar{C_{in}} + ABC_{in}
\]

\textbf{Look at that.} It is $\sum m(3,5,6,7)$ --- \emph{the same minterm set as
Exercise 1}, with $C_{in}$ playing the role of $C$. The derivation is therefore the
one you already did, and the answer is:
\[
C_{out} = \boxed{AB + AC_{in} + BC_{in}}
\]

\textbf{$C_{out}$ is the majority function.} And it must be: a carry is produced
exactly when at least two of the three input bits are $1$. The first circuit you
built this semester is the carry logic of an adder.

\textbf{(3) VHDL.}
\begin{lstlisting}
Cout <= (A and B) or (A and Cin) or (B and Cin);
\end{lstlisting}

\subsection*{4.5 The structural architecture}

The same entity, built from two half adders and an OR gate.

Half adder: $S = A \oplus B$, \quad $C = A \cdot B$.

\begin{lstlisting}
architecture structural of full_adder is
    signal S1, C1, C2 : std_logic;
begin
    HA1 : entity work.half_adder
        port map ( A => A,  B => B,   S => S1,  C => C1 );
    HA2 : entity work.half_adder
        port map ( A => S1, B => Cin, S => Sum, C => C2 );
    Cout <= C1 or C2;
end architecture structural;
\end{lstlisting}

Note \textbf{named association} in both \code{port map}s --- \code{A => A} rather
than positional order. Use it everywhere, all semester. Positional port maps are the
single most common source of silently-wrong structural designs, and they get worse
as components grow: the register file in Lab 6 has seven ports.

\subsection*{4.6 Proving the two architectures agree}

The structural version produces
\[
C_{out} = C_1 + C_2 = AB + (A \oplus B)C_{in}
\]
which does not \emph{look} like $AB + AC_{in} + BC_{in}$. Expand the behavioural
form to see that it is:
\begin{align*}
AB + AC_{in} + BC_{in}
&= AB + AC_{in}(B + \bar{B}) + BC_{in}(A + \bar{A}) && \text{\small complement law, backwards} \\
&= AB + ABC_{in} + A\bar{B}C_{in} + ABC_{in} + \bar{A}BC_{in} \\
&= AB(1 + C_{in} + C_{in}) + A\bar{B}C_{in} + \bar{A}BC_{in} && \text{\small distributive law} \\
&= AB + A\bar{B}C_{in} + \bar{A}BC_{in} && \text{\small annulment: } 1 + X = 1\\
&= AB + (A\bar{B} + \bar{A}B)C_{in} \\
&= AB + (A \oplus B)\,C_{in} \qquad\checkmark
\end{align*}

Equal. The structural version reuses the $A \oplus B$ term it already needed for the
sum, so it costs one fewer gate --- which is why real adders are built this way.

\textbf{But do not take the proof's word for it.} Instantiate both architectures in
one testbench, drive both from the same stimulus, and assert they agree on all eight
vectors:

\begin{lstlisting}
assert (Sum_beh = Sum_str) and (Cout_beh = Cout_str)
    report "row " & integer'image(i) & ": architectures disagree"
    severity error;
\end{lstlisting}

That one line is a habit worth forming now. It is how you will check every dual
implementation for the rest of the course.

\subsection*{4.7 Where this goes}

Chain 32 full adders carry-to-carry and you have the 32-bit adder inside the ALU you
build in Lab 5. The carry out of bit 31 is what tells you the result overflowed.

\vspace{6pt}\hr

\section*{Summary --- the four results}

\begin{center}
\begin{tabular}{clll}
\toprule
& circuit & result & reappears in \\
\midrule
1 & majority-of-3 & $Y = AB + AC + BC$ & Ex.\ 4, as $C_{out}$ \\
2 & 2:1 mux & $Y = \bar{S}D_0 + SD_1$ & Lab 3, 32 bits wide \\
3 & 2-to-4 decoder & $D_k = E \cdot m_k$ & Lab 6, register file writes \\
4 & full adder & $Sum = A \oplus B \oplus C_{in}$ & Lab 5, ALU \\
  &             & $C_{out} = AB + AC_{in} + BC_{in}$ & \\
\bottomrule
\end{tabular}
\end{center}

Four circuits, and three of them are inside the processor you finish in December.

\vspace{4pt}
\hr
{\small \textbf{Three things to carry forward.} A K-map checkerboard means reach for
XOR. Named association in every \code{port map}. And open the RTL Viewer before you
compile --- Exercise 3's bug was visible there in ten seconds.}

\end{document}
